% Response to reviewers -- ICMLDE 2026
% Companion to the revised manuscript "Why AI-Generated Image Detectors Fail
% on Modern Diffusion Generators: A Diagnostic Study of Cross-Architecture
% Generalization and Failure Modes".
% Section, table and figure numbers below refer to the revised manuscript.
\documentclass[11pt]{article}
\usepackage[a4paper,margin=1in]{geometry}
\usepackage{times}
\usepackage[T1]{fontenc}
\usepackage{microtype}
\usepackage{url}
\setlength{\parskip}{3pt}
\setlength{\parindent}{0pt}

\title{\vspace{-2.2em}Response to Reviewers' Comments}
\author{}
\date{}

\begin{document}
\maketitle
\vspace{-3.2em}

\noindent\textbf{Paper:} Why AI-Generated Image Detectors Fail on Modern
Diffusion Generators: A Diagnostic Study of Cross-Architecture
Generalization and Failure Modes\\
\textbf{Authors:} Kavinesh P, Raam Prathap R V, Santhosh A S, Kasam Sai
Venkata Siddardha, Ashish M, Anbazhagan Mahadevan\\
\textbf{Conference:} 5th International Conference on Machine Learning and
Data Engineering (ICMLDE 2026)

\vspace{6pt}
\noindent All section, table and figure numbers below refer to the revised
manuscript.

\vspace{4pt}
\hrule
\vspace{10pt}

We thank the three reviewers for comments that were specific enough to act
on directly. Every point is answered below with the location of the
corresponding change. Four of the suggestions required new experiments,
all of which we ran: a test of the ensembling hypothesis
(Section~4.4), per-detector reliability diagrams
(Figure~1), a matched-compression condition that
quantifies the PNG/JPEG confound (Section~6), and a
comparison against published numbers for the same checkpoints
(Section~3.5). One suggestion, adding a further
contemporary detector, we were only able to address in part, and we say so
plainly below rather than claiming otherwise.

\section*{Reviewer 1}

\noindent\textbf{1.1 Abstract is compacted to highlight the main findings.}\\[1pt]
The abstract has been rewritten. It now opens by stating that this is a
diagnostic study, gives the per-generator sample sizes, and reports the
new ensembling and matched-compression results, while the preprocessing
and compression cautions are stated as bounds on the absolute numbers
rather than as closing remarks.

\noindent\textbf{1.2 The systematic analysis of failure profiles and
calibration under distribution shift is a major contribution.}\\[1pt]
Thank you. We have strengthened this part rather than leaving it as it
was: Section~4.7 now carries per-detector reliability
diagrams (Figure~1, also requested by Reviewer 3),
which show that for every failing cell the largest confidence bin is the
rightmost one. That makes the point visible in a way the scalar ECE in
Table~4 cannot: the curves fall vertically rather than
shifting toward lower confidence.

\noindent\textbf{1.3 Training data finding is based on two detectors and
should be appropriately qualified.}\\[1pt]
Agreed, and we have done two things. First, the closing paragraph of
Section~4.3 now marks the claim as a hypothesis rather
than a finding and states the specific weaknesses: it rests on one
controlled pair, $n=5$ generators, the permutation $p$ of 0.017 is the
floor of what the design can produce rather than an effect size, and a
shared quirk of the SDv1.4 training subset could produce the same ordering
twice for unrelated reasons. The abstract and conclusion were softened to
match. Second, rather than only qualifying it, we tested a falsifiable
prediction it makes (Section~4.4); the hypothesis survived
that test, which is better evidence than restating it more cautiously
would have been.

\noindent\textbf{1.4 The leakage verification and McNemar tests are solid
and should be kept as examples of evaluation practice.}\\[1pt]
Both are retained in full (Sections~3.4
and~4.2, Table~5). The McNemar
machinery is also now reused for the new ensembling comparison.

\noindent\textbf{1.5 Including at least one contemporary detector for
diffusion models would widen the comparison.}\\[1pt]
Addressed in part, and we want to be accurate about which part. The
evaluated set already contains two detectors built for the diffusion era,
NPR (CVPR 2024) and the OMAT-hardened CLIP-LoRA (2025), the latter
designed and published specifically against modern diffusion and Diffusion
Transformer outputs; Table~1 now carries a Year column
so this span is visible rather than having to be inferred. What we could
not do within the revision window is add a detector from a structurally
independent family, meaning one whose detection signal is neither a CLIP
embedding nor a pixel residual. We went further than asserting this. The
obvious candidate is DIRE, so we obtained its released checkpoint, built a
single-GPU harness for it, and measured the cost: DIRE scores an image by
inverting it to a latent with DDIM and sampling back, roughly 40 diffusion
forward passes, which ran at 171 seconds per image on the 8~GB GPU used
throughout. That puts 100 images per class on one generator at 9.5
GPU-hours, and the 6{,}000-per-class protocol the other five detectors run
on, across five generators, near 2{,}800. Reporting a handful of images
instead would have been the partial sweep we said we would not report, so
Section~8 now carries the measurement and the harness ships with the code,
which means the experiment is reproducible on larger hardware rather than
merely deferred.
Section~8 states this as a limitation and names it as
the most valuable single addition future work could make, rather than
leaving the reader to discover the gap.

\noindent\textbf{1.6 Adding a short list of practical recommendations for
practitioners would be useful.}\\[1pt]
Added as Section~7: six recommendations for
selecting and deploying a detector, each tied to a specific measurement in
the paper, plus one recommendation addressed to benchmark authors
regarding scriptable verification subsets.

\section*{Reviewer 2}

\noindent\textbf{2.1 Confounding due to JPEG vs PNG compression is an
important note that could be highlighted earlier in the abstract.}\\[1pt]
Moved. The compression confound is now stated in the abstract with its
measured size, and it appears in the introduction's contribution list
rather than only at the end of the paper. It is also no longer a
subsection buried inside a robustness section: Section~6
treats the robustness question and the benchmark confound together,
because the single experiment answers both.

\noindent\textbf{2.2 The analysis of resize conventions is a practical
contribution that deserves to be highlighted.}\\[1pt]
Promoted. It is now one of the five numbered contributions in the
introduction, has its own top-level section
(Section~5) rather than sitting under a discussion
heading, and generates recommendation~4 in
Section~7. The result that two checkpoints trained
on identical data prefer opposite preprocessing is stated in the abstract.

\noindent\textbf{2.3 The manuscript should clearly state its nature as a
diagnostic study.}\\[1pt]
Done in three places. The title now reads ``A Diagnostic Study of
Cross-Architecture Generalization and Failure Modes''; the abstract's first
sentence says the study trains nothing and proposes no new detector; and
the introduction has a short paragraph making the same point before the
contribution list.

\noindent\textbf{2.4 Being over 8,000 words long, the manuscript is too
long for a conference and should be shortened.}\\[1pt]
Material carried over from the submitted version has been cut by roughly a
quarter, and the revised manuscript is shorter than the one the reviewers
read despite absorbing four new experiments, three new tables and a new
figure. The cuts came from merging the three preprocessing subsections
into one section with a single subsection; reducing the spectral-deviation
analysis from roughly a page to a short paragraph and removing its figure,
since it is a negative result that did not need the space; folding the
``generator easiness is detector-relative'' subsection into Section~4.1;
removing the per-image score histogram, whose evidence the new reliability
diagrams carry more rigorously; condensing the published-number comparison
in Section~3.5; rewriting the conclusion so that it draws out consequences
instead of restating the abstract; and removing the numbers the abstract,
introduction and results had been repeating to each other. No result was
dropped. This response is submitted
alongside the manuscript and is not part of the manuscript length.

\noindent\textbf{2.5 Results for matched compression in real and synthetic
cases would give a quantitative estimate of the confound effect.}\\[1pt]
This was the most useful suggestion we received and we ran the experiment.
Section~6 and Table~8 now report
three conditions on one fixed FLUX.1 sample of 2{,}000 images per class,
paired within each detector: the benchmark as shipped; a
\emph{fakes-only} condition that re-encodes generated images at JPEG q75
and leaves real images untouched, which removes the format asymmetry and
nothing else; and a condition re-encoding everything. The result is
sharper than we expected. Removing only the format asymmetry takes NPR
from AUC 97.10 to 39.13, which is below chance, meaning the detector then
ranks real photographs as more likely generated than the generated images.
We note in the text why that matters for interpretation: destroyed
evidence would drive AUC toward 50 and stop, so an inversion indicates the
quantity NPR reads tracks relative compression level rather than
provenance. The same manipulation costs the CLIP baseline 4.5 AUC points,
the ResNet-50 4.3, CLIP-LoRA 12.8 and UniversalFakeDetect 1.0, against
58.0 for NPR, so the confound inflates an artifact-based detector by an
order of magnitude more than a semantics-based one. For comparison,
Grommelt et al.\ report shifts above 11 points when removing JPEG and
dimension confounds from GenImage.

\noindent\textbf{2.6 The ensembling hypothesis could be tested in a small
experiment to support the conclusion.}\\[1pt]
Done, as a new Section~4.4 with Table~6.
The submitted version predicted that combining detectors chosen for
different training lineages should beat combining for architectural
variety; the per-image scores already collected let us test that without
retraining. Holding the GenImage++ CLIP baseline fixed and varying only
its partner, a same-lineage partner of different architecture takes
worst-generator fake-accuracy from 9.6\% down to 0.3\%, while a
different-lineage partner lifts it to 62.9\%. A paired McNemar test
separates the two combinations on all five generators
($p<10^{-200}$ throughout). Across all pairs, excluding the nested
CLIP/CLIP-LoRA pair, cross-lineage pairs average 52.5\% worst-generator
accuracy against 37.3\% for same-lineage pairs. We also report two
findings that cut against a simple ``ensemble more'' message and would
have been easy to omit: lineage diversity is necessary but not sufficient,
since a cross-lineage pair with a weak partner underperforms the strong
detector alone, and every pair we construct is beaten by the single best
detector. The conclusion is phrased accordingly, as a statement about
where blind spots come from rather than a leaderboard claim.

\section*{Reviewer 3}

\noindent\textbf{3.1 Abstract should mention the sample sizes for each
generator.}\\[1pt]
Added: 6{,}000 real and 6{,}000 generated images for BigGAN, GLIDE, FLUX.1
and SD3, 18{,}300 per class for SDXL, and 2{,}000 per class for the paired
compression experiments.

\noindent\textbf{3.2 Sharing the evaluation code and preprocessing
parameters would improve the significance of the work.}\\[1pt]
We agree, particularly given that Section~5 argues
precisely this point about other people's releases. Appendix~A now
itemizes what is released rather than promising a repository in general
terms: the harness and the exact command line for each of the 25
detector/generator cells, the per-image prediction scores every table and
test in the paper is computed from, the preprocessing functions for the
three resize conditions together with the normalization statistics used
per detector, and the scripts that regenerate every results table and the
figure from the score files. The repository is public at
\url{https://github.com/MogSquad-FinalYear/BlindSpot}.

\noindent\textbf{3.3 Calibration analysis could be extended by adding
reliability diagrams for each detector.}\\[1pt]
Added as Figure~1: one panel per detector, five
generators each, plotted in the same coordinates the reported ECE
integrates so the figure and Table~4 are consistent by
construction. Marker area is proportional to bin occupancy, which we found
necessary because almost all mass sits at confidence near 1 and an
unweighted curve makes sparse middle bins look as informative as the
extremes; bins with fewer than 25 images are omitted. The final paragraph
of Section~4.7 reports what the diagrams add beyond the
scalar.

\noindent\textbf{3.4 Choosing the five generators should be motivated in
terms of covering current generator families.}\\[1pt]
Section~3.2 now opens with that motivation. The set
covers the three denoiser and generator families that account for
essentially all current open-weight synthesis, with the U-Net family given
two members (GLIDE and SDXL) so it contributes a within-family contrast of
four years rather than a single point. We also name the two families the
study does not cover, autoregressive token-based generators and
face-specific manipulation pipelines, the second deliberately.

\noindent\textbf{3.5 Comparisons with results from GenImage and other
benchmarks would help the reader to relate these findings.}\\[1pt]
Added as Section~3.5 and Table~2,
placed in the methodology so the harness is validated before new numbers
are reported. We compare our overall accuracy against the figures the NPR
repository publishes for the same ProGAN-4class checkpoint we evaluate:
three of four cells agree within five points and the UniversalFakeDetect
GLIDE cell agrees to within 0.1. We flag that those subsets are assembled
differently from GenImage's, so this is a consistency check rather than a
reproduction. The one same-benchmark comparison available, against Zhou et
al.'s reported 62.66\% for their CLIP baseline on BigGAN, is the
disagreement that Section~5 resolves, and
Section~3.5 now forward-references it instead of
leaving the reader to encounter the 6x gap unprepared.

\end{document}

\end{document}
